% Layout adapted from the instructor's ME 537 homework sources.
\documentclass[12pt,letterpaper]{article}
\usepackage{mathptmx}
\usepackage[T1]{fontenc}
\usepackage{xcolor,amsmath,amssymb,xurl,listings,enumitem}
\textwidth=6.5in
\textheight=9in
\oddsidemargin=0in
\evensidemargin=0in
\topmargin=0in
\headheight=0in
\headsep=0in
\parskip=0.10in
\parindent=0in
\setlist[enumerate]{itemsep=0.05in,topsep=0.04in}
\lstset{basicstyle=\ttfamily\footnotesize,breaklines=true,columns=fullflexible,
  keepspaces=true,showstringspaces=false,aboveskip=0.05in,belowskip=0.05in}

\begin{document}
\setlength{\fboxrule}{2pt}\setlength{\fboxsep}{2mm}
\fcolorbox{black}{white}{\parbox{\dimexpr\linewidth-2\fboxsep-2\fboxrule\relax}{\centering\textbf{Introduction to Robotics Homework \#1 -- Draft}}}

\textbf{ROS 2 setup, messages, and motion over time}

This assignment checks your environment and introduces a small application:
publishing robot joint positions. You do not need prior robotics experience
or forward kinematics. The supplied robot is a visualization, not a physics
simulation. Use the student repository supplied by the instructor.

\begin{enumerate}
\item \textbf{Install, start, and see the robot.}
Read \path{docs/setup.md} and the guide for your operating system.
\begin{enumerate}
\item Install and start Docker Desktop on Windows or macOS. Windows uses
Linux containers with WSL 2. Linux users follow the Docker Engine directions
in the setup guide. Open the course repository folder in your editor.
\item In a terminal \emph{on your computer}, change to that repository
folder (the one containing \path{compose.yaml}) and run:
\begin{lstlisting}
docker compose build
docker compose up -d
\end{lstlisting}
\item Open your usual web browser and visit this exact address:

\url{http://localhost:6080/vnc.html?autoconnect=1&resize=scale}

The page shows the container's Linux desktop and RViz. Docker Desktop's
dashboard is a different application. Wait for the first robot-package build
to finish. If needed, inspect \texttt{docker compose logs -f course}; Ctrl+C
stops following the logs without stopping the container.
\item Find the Joint State Publisher window. It may be \emph{behind RViz};
move or minimize RViz to reveal it. Move the head and both leg sliders.
Select \textbf{Move Camera} in RViz and drag in the 3D view; scroll to zoom.
Save a screenshot with the robot and sliders visible and the head turned.
\item Record your host operating system, Docker version, and image
architecture. Run these in the host terminal:
\begin{lstlisting}
docker version
docker image inspect intro-robotics-course:jazzy --format '{{.Os}}/{{.Architecture}}'
\end{lstlisting}
An Apple Silicon image should be \texttt{linux/arm64}; an x64 Windows image
should be \texttt{linux/amd64}. Report setup difficulties and resolutions.
\end{enumerate}
\end{enumerate}

\newpage
\begin{enumerate}[start=2]
\item \textbf{Publish a fixed joint pose from Python.}
Disable the automatic slider publisher before running your node. In a
\emph{host} terminal in the repository folder, run:
\begin{lstlisting}
docker compose -f compose.yaml -f compose.publisher.yaml up -d --force-recreate
docker compose exec course bash
\end{lstlisting}
Refresh the browser page after the container is recreated. The second
command opens a Linux shell \emph{inside the container}; ROS 2 is available
there. Robot links may not appear until your first valid joint-state message;
missing-transform messages before that are expected. In that shell, run:
\begin{lstlisting}
cd /workspace
python3 homework/hw01/joint_publisher.py
\end{lstlisting}
The starter intentionally raises a homework-labeled
\texttt{NotImplementedError}. Edit it using your host editor. Stop the node
with Ctrl+C and rerun after edits; this Python file needs no image rebuild.

The supplied \texttt{JointPublisher} class owns a publisher, a start time,
and a timer. Complete these two methods:
\begin{enumerate}
\item \texttt{joint\_positions(elapsed)}: initially return a fixed list
of three joint angles in radians, in the supplied joint-name order. Choose
a visibly changed head pose and keep both legs within their limits.
\item \texttt{publish\_joint\_state()}: create a \texttt{JointState}
message, fill its timestamp, joint names, and positions, and publish it with
\texttt{self.publisher}. The provided variables are \texttt{now},
\texttt{positions}, and \texttt{JOINT\_NAMES}. For the timestamp use the
ROS message form of \texttt{now}; consult the linked ROS tutorial.
\end{enumerate}
\texttt{support.py} supplies names and the explained
\texttt{check\_joint\_positions} utility. You do not need to implement it.

\begin{center}
\begin{tabular}{lll}
\hline
Position in list & Joint name & Allowed angle (radians) \\
\hline
1 & \texttt{head\_yaw} & $[-1.57,\ 1.57]$ \\
2 & \texttt{left\_leg\_tilt} & $[-0.5,\ 0.5]$ \\
3 & \texttt{right\_leg\_tilt} & $[-0.5,\ 0.5]$ \\
\hline
\end{tabular}
\end{center}

\item \textbf{Make the motion depend on time.}
The provided timer calls your publishing method every 0.1 seconds.
\begin{enumerate}
\item Confirm that repeatedly publishing a fixed pose holds the same pose.
Then change \texttt{joint\_positions(elapsed)} so the head moves smoothly
back and forth with a period between 4 and 10 seconds. Keep all angles
within their limits. Moving the legs smoothly is an optional extension.
\item Compute positions from elapsed time; keep the timer callback short.
Do not add a blocking \texttt{while} loop or a sleep inside that callback.
Save a short screen recording, or two time-labeled screenshots with message
evidence, showing different poses while the same node continues running.
\end{enumerate}
\end{enumerate}

\newpage
\begin{enumerate}[start=4]
\item \textbf{Inspect messages and explain what is happening.}
While your node runs, open a second host terminal in the repository folder
and enter a second container shell with \texttt{docker compose exec course bash}.
Run the following \emph{inside} that second shell:
\begin{lstlisting}
ros2 topic echo /joint_states --once
ros2 topic hz /joint_states
ros2 topic info /joint_states --verbose
\end{lstlisting}
Stop the rate display with Ctrl+C. Save one message and the measured rate
(about 10 Hz; report what you measure), and check that your node is the only
publisher of \texttt{/joint\_states}. Answer briefly:
\begin{enumerate}
\item How do the joint names correspond to the values in \texttt{position}?
What are the position units? What does the message timestamp represent?
\item Why does repeated publishing of an unchanged position not produce
motion? How does your time-dependent function change that?
\item What does the timer do, and what does \texttt{rclpy.spin(node)} do?
Describe one use of \texttt{self} in the provided class.
\item Explain this chain in your own words:
\begin{center}\small
\texttt{/joint\_states} $\longrightarrow$
\texttt{robot\_state\_publisher} $\longrightarrow$
\texttt{/tf} $\longrightarrow$ RViz.
\end{center}
\item Plot your motion over time. In the second shell, run:
\begin{lstlisting}
ros2 run plotjuggler plotjuggler
\end{lstlisting}
Set \textbf{Buffer} to 30~s, choose \textbf{ROS2 Topic Subscriber} under Streaming, click \textbf{Start},
select \texttt{/joint\_states}, and click OK. Drag
\texttt{/joint\_states/head\_yaw/position} into the plot. Save a screenshot
of at least two periods; compare the period and amplitude you read with your
\texttt{joint\_positions} function.
\end{enumerate}

\item \textbf{Stop and restart safely.}
Stop your Python node with Ctrl+C. In a \emph{host} terminal, stop the course
container with:
\begin{lstlisting}
docker compose down
\end{lstlisting}
Closing only the browser tab does not stop the container; your files remain
in the repository folder. To return to the original sliders later, start Docker
Desktop and run \texttt{docker compose up -d -{}-force-recreate}, then reopen
the same web address.
\end{enumerate}

\textbf{Submit:} your completed Python file; setup screenshot and environment
record; fixed-pose and timed-motion evidence; one joint-state message and
measured rate; the PlotJuggler plot of the head angle; the short explanations
above; and a statement of AI use
(including ``none'' if applicable). We will review correct behavior,
readable code, and your understanding of the provided tools. Follow the instructor's submission destination and deadline.

\textbf{References and troubleshooting:} copyable commands and
troubleshooting notes are in \path{docs/setup.md} and
\path{homework/hw01/README.md}. ROS publisher tutorial:
\url{https://docs.ros.org/en/jazzy/Tutorials/Beginner-Client-Libraries/Writing-A-Simple-Py-Publisher-And-Subscriber.html}.
Message fields: \url{https://docs.ros.org/en/jazzy/p/sensor_msgs/msg/JointState.html}.
\end{document}
